Everything is computed in Gauge Theory Gravity (Definition \claimref{obj.synthesis.gtg-flat-spacetime-gauge-fields}) and in the conformal model. Physical vectors at $x$ use the flat product, so $g(v, w) = v\cdot w$; the Schouten tensor is the symmetric linear map $\Sch(a) = \tfrac12(\mathcal{R}(a) - a\mathcal{R}/6)$ of the field strength, with $\Sch(v, w) = v\cdot\Sch(w)$ and $J = \partial_{a}\cdot\Sch(a)$ (Fact \claimref{synthesis.imported-gtg-dictionary}); $\Ric(v, w) = v\cdot\mathcal{R}(w)$ in the store's sign, and $T(v, w) = v\cdot T(w)$ for a symmetric linear map $T$.

\emph{Setting.} A tractor at $x$ is a vector $\sigma\nInf + \mu + \rho\,\nO$ of $\Cl(2,4)$ with $\mu$ a physical vector and $\nO\cdot\nInf = 1$, so the tractor form is the flat product, $(\sigma\nInf + \mu + \rho\nO)\cdot(\sigma'\nInf + \mu' + \rho'\nO) = \sigma\rho' + \rho\sigma' + \mu\cdot\mu'$, and $I_g = \nInf$, $X = \nO$. By Theorem \claimref{synthesis.gtg-tractor-gauge-field-reproduces-tractor-connection}, the connection of Fact \claimref{interior.imported-tractor-connection} is the gauge field $\Omega_{T}(\hb(v)) = \Omega(\hb(v)) - \nInf\wedge v + \nO\wedge\Sch(v)$ along a physical vector $v$, acting by $D_{v}V = (\hb(v)\cdot\nabla)V + \Omega_{T}(\hb(v))\times V$. A bivector acting by the commutator preserves the flat product, which is Lemma \claimref{interior.metric-parallelism-trace-schouten-condition}(i). On the constant vectors $\nInf$ and $\nO$, the spacetime bivector $\Omega(\hb(v))$ commutes with both, and with $(a\wedge b)\cdot c = a\,(b\cdot c) - b\,(a\cdot c)$:
\[
D_{v}\nInf = (\nO\wedge\Sch(v))\cdot\nInf = -\Sch(v), \qquad D_{v}\nO = -(\nInf\wedge v)\cdot\nO = v,
\]
the slot forms $(0, -\Sch(v,\cdot), 0)$ and $(0, v^{\flat}, 0)$; the first is Theorem \claimref{interior.covariant-form-vacuum-equations}. The check computes $\Omega_{T}$ on five gauge fields (de Sitter in the flat slicing, Schwarzschild-de Sitter and Schwarzschild in the Painlev\'e-Gullstrand chart, a flat FLRW instant that is not Einstein, a Kasner vacuum), with $\Sch$ from each field strength, and finds these two derivatives along every frame vector and the flat product preserved.

\emph{Complex numbers.} The complexified tractor space is the span of $V + I_{7}V'$ in $\Cl(3,4) = \Cl(2,4) \oplus \Cl(2,4)I_{6}e'$, with $Q(e') = -I_{6}^{2} = 1$ and $I_{7} = I_{6}e'$ central of square $-1$ (the unit principle, U3; \claimref{std-algebra.complexification-is-the-odd-extension}); its numbers are $\mathrm{span}\{1, I_{7}\}$, the form extended bilinearly is the symmetric product $\tfrac12(xy + yx)$, and the $\mu$ slot of $\mathcal{T}_x\otimes\C$ is the complexified tangent space with the same unit. Conjugation $\bar{\ }$ changes the sign of the $I_{7}$ part.

\emph{(1) Shape.} Since $V\cdot\nInf = \rho$, the orthogonal complement of $I_g$ is $\{\rho = 0\}$, and a totally isotropic $W \ni I_g$ lies in it. Subtracting multiples of $I_g$ from two further basis vectors removes their $\nInf$ components, so $W = \mathrm{span}(I_g, \ell, m)$ with $\ell, m$ in the $\mu$ slot, and isotropy of $W$ is exactly total nullity of $A_W = \mathrm{span}(\ell, m)$. The same argument with $V\cdot\nO = \sigma$ handles point vacua. A totally null plane $A$ of the four-dimensional complexified tangent space equals its own orthogonal complement. If $A \cap \bar A = 0$, every real vector is $a + \bar a$ with $a \in A$, and $Q(a + \bar a) = 2\,a\cdot\bar a$, so the Hermitian form $h(a, b) = a\cdot\bar b$ on the two-dimensional $A$ would give the real space a signature $(2p, 2q)$, contradicting $(1,3)$. So $A \cap \bar A \neq 0$; it is conjugation-invariant, so it contains a nonzero real vector $\ell_W$, null since $A$ is totally null. Two real null vectors in $A$ are orthogonal, and orthogonal real null vectors in signature $(1,3)$ are proportional, so the real null direction is unique. Conversely, for a real null $\ell$ the quotient $\ell^{\perp}/\ell$ is two-dimensional with definite induced form, whose complexification has exactly two null lines; their preimages are the two totally null planes containing $\ell$ (for $\ell = e_{0} + n$, $n$ a unit spatial vector and $s_{1}, s_{2}$ an orthonormal basis of the screen, they are $\mathrm{span}(\ell, s_{1} \pm I_{7}s_{2})$), so $\ell$ lies in exactly two infinity vacua. The identification with Fock vacua is the purity statement of Fact \claimref{one-relation.imported-pure-spinors-incidence-chart}: each conformal spinor $u$ has annihilator $W_u = \{v : c(v)u = 0\}$ of dimension three, determining $u$ up to scale, so for $u \in \CS$, $I_g \in W_u$ if and only if $c(I_g)u = 0$. The kernel of $c(I_g)$ is half the eight-dimensional Fock module: $c(I_g)^{2} = Q(I_g) = 0$ gives $\mathrm{im}\,c(I_g) \subseteq \ker c(I_g)$, and $c(I_g)c(X) + c(X)c(I_g) = 2\,I_g\cdot X = 2$ gives $\psi = \tfrac12 c(I_g)c(X)\psi \in \mathrm{im}\,c(I_g)$ for $\psi \in \ker c(I_g)$, so kernel and image coincide. The two totally null planes through a real null direction belong to the two families (two maximal isotropic subspaces of the six-dimensional complexified space lie in the same family exactly when their intersection has odd dimension, and the two infinity vacua through $\ell$ meet in $\mathrm{span}(I_g, \ell)$).

\emph{(2) Classes.} For $w' \in W$ and a function $f$, $(D_{v}(fw))\cdot w' = f\,(D_{v}w)\cdot w' + (vf)\,w\cdot w' = f\,(D_{v}w)\cdot w'$ because $W$ is isotropic, so $\omega_{v}$ is tensorial in $w$, and by $(D_{v}w)\cdot w' = -w\cdot D_{v}w'$ in $w'$ as well; it depends only on $W$ at $x$. With the two derivatives of the setting, for $w'$ in the $\mu$ slot: $\omega_{v}(I_g, w') = -\Sch(v)\cdot w' = -\Sch(v, w')$ and $\omega_{v}(X, w') = v\cdot w' = g(v, w')$, while $\omega_{v}(I_g, I_g) = \omega_{v}(X, X) = 0$. For the Fock statement (with $W$ in the family of $\CS$), $c(D_{v}I_g)u_W = 0$ if and only if $D_{v}I_g \in W_{u_W} = W$; a $\mu$-slot vector lies in $W$ exactly when it lies in $A_W$, and since $A_W = A_W^{\perp}$, exactly when it is orthogonal to $\ell$ and $m$, that is when $\omega_{v}(I_g, \ell) = \omega_{v}(I_g, m) = 0$.

\emph{(3) Point vacua.} $\omega_{\ell_W}(X, w') = \ell_W\cdot w'$ vanishes on $A_W$ because $A_W$ is totally null and contains $\ell_W$; no field strength enters.

\emph{(4) Einstein.} (b) implies (a) trivially. For (a) implies (c): by (1) every real null direction is some $\ell_W$, so (a) says $\ell\cdot\Sch(\ell) = 0$ for every real null $\ell$. A symmetric linear map $S$ with $\ell\cdot S(\ell) = 0$ on the null cone is a multiple of the identity: in an orthonormal frame, $\ell = e_{0} \pm e_{i}$ give $S_{00} + S_{ii} = 0$ and $S_{0i} = 0$ (with $S_{ab} = e_{a}\cdot S(e_{b})$), and $\ell = e_{0} + (\cos\theta)e_{i} + (\sin\theta)e_{j}$ with $\cos\theta, \sin\theta$ both nonzero then gives $S_{ij} = 0$ for $i \neq j$; so $S(a) = c\,a$ with $c = S_{00}$. Hence $\Sch(a) = c\,a$, and the trace gives $c = J/4$; Lemma \claimref{interior.metric-parallelism-trace-schouten-condition}(ii) gives $c = \mathcal{R}/24$ and $\mathcal{R}(a) = \tfrac{\mathcal{R}}{4}a$, that is $\Ric = \tfrac{R}{4}g$. For (c) implies (b): if $\Sch(a) = \tfrac{J}{4}a$ then $\omega_{\ell_W}(I_g, w) = -\tfrac{J}{4}\,\ell_W\cdot w = 0$ for $w \in A_W$ by total nullity. Constancy of $R$ on a connected region where $\mathcal{R}(a) = \tfrac{\mathcal{R}}{4}a$: the contracted Bianchi identity of the field strength, $\dot{\mathcal{R}}(\dot D) = \tfrac12 D\mathcal{R}$ (the covariant conservation of the Einstein map, $\dot{\mathcal{G}}(\dot D) = 0$; an identity of the torsion-free gauge field, unchanged by the overall sign convention since both sides change sign together), gives $\tfrac14 D\mathcal{R} = \tfrac12 D\mathcal{R}$, so the physical gradient $D\mathcal{R} = \hb^{T}(\nabla\mathcal{R})$ vanishes and $\mathcal{R}$ is constant on the connected region (Schur's argument in dimension four); then $J = \mathcal{R}/6$, $\Sch(a) = \tfrac{\Lambda}{6}a$ with $\Lambda = \mathcal{R}/4$, and $\omega_{v}(I_g, w) = -\tfrac{\Lambda}{6}\,v\cdot w$ by (2).

\emph{(5) Ricci-flat.} If every infinity vacuum has vanishing $(I_g, \cdot)$ classes along every $v$, then $v\cdot\Sch(\ell) = 0$ for every $v$ and every real null $\ell$, and real null vectors span the physical vectors, so $\Sch = 0$; conversely $\Sch = 0$ makes every class $-\Sch(v, \cdot)$ vanish. $\Sch = 0$ if and only if $\mathcal{R}(a) = 0$ by Lemma \claimref{interior.metric-parallelism-trace-schouten-condition}(ii), and, by the setting, if and only if $D_{v}I_g = 0$ for every $v$ (Theorem \claimref{interior.covariant-form-vacuum-equations}).

\emph{(6) One vacuum.} The condition is $v\cdot\Sch(a) = 0$ for all $v$ and all $a \in A = A_W$, that is $\Sch(a) = 0$ on $A$ (extended linearly over $\mathrm{span}\{1, I_{7}\}$). Since $\Sch$ is real it also vanishes on $\bar A$, hence on $A + \bar A$, which has dimension three because $A \cap \bar A$ is the line of $\ell_W$ by (1); so $A + \bar A$ is the complexification of $\ell_W^{\perp}$. By symmetry the image of $\Sch$ is orthogonal to $\ell_W^{\perp}$, so it lies in the line of $\ell_W$, and $\Sch(a) = c\,(a\cdot\ell_W)\,\ell_W$ with $c$ real (the map is real and symmetric). This is $\Sch = c\,\ell_W^{\flat}\otimes\ell_W^{\flat}$. Conversely such a map has $v\cdot\Sch(a) = c\,(v\cdot\ell_W)(\ell_W\cdot a) = 0$ for $a \in A$. Its trace is $c\,Q(\ell_W) = 0$, so $J = \mathcal{R}/6 = 0$ and $\mathcal{R}(a) = 2\Sch(a) + \tfrac{\mathcal{R}}{6}a = 2c\,(a\cdot\ell_W)\,\ell_W$.

\emph{(7) Matter.} For null $\ell$, $\ell\cdot\Sch(\ell) = \tfrac12\,\ell\cdot\mathcal{R}(\ell)$, so the condition $-\ell\cdot\Sch(\ell) = 4\pi G\,\ell\cdot T(\ell)$ reads $\ell\cdot(\mathcal{R} + 8\pi G\,T)(\ell) = 0$ for every real null $\ell$, and the argument of (4) applied to the real symmetric map $\mathcal{R}(a) + 8\pi G\,T(a)$ gives $\mathcal{R}(a) + 8\pi G\,T(a) = h\,a$; the converse is immediate. Now let this hold on a connected region with $\dot T(\dot D) = 0$ (the conservation $\nabla^{a}T_{ab} = 0$), and write $\tau := \partial_{a}\cdot T(a)$. The trace gives $\mathcal{R} + 8\pi G\,\tau = 4h$. The divergence gives $\dot{\mathcal{R}}(\dot D) = Dh$, and the contracted Bianchi identity gives $\dot{\mathcal{R}}(\dot D) = \tfrac12 D\mathcal{R} = \tfrac12 D(4h - 8\pi G\tau) = 2Dh - 4\pi G\,D\tau$. Hence $Dh = 4\pi G\,D\tau$, and $\Lambda := h - 4\pi G\tau$ is constant. For the sign remark: the store's curvature has $\mathcal{R}(a) = \Lambda a$ on de Sitter, and the round-sphere-positive Ricci tensor is $\Ric^{\circ}(u, v) = -u\cdot\mathcal{R}(v)$ (dictionary C6, C11); the Einstein equations $\Ric^{\circ} = 8\pi G(T - \tfrac12\tau g) - \Lambda g$ in signature $(1,3)$ then read $\mathcal{R}(a) + 8\pi G\,T(a) = (\Lambda + 4\pi G\,\tau)\,a$, the GTG field equation $\mathcal{G}(a) + \Lambda a = -8\pi G\,T(a)$ of the dictionary, which is the form stated.

\emph{The checks.} The GTG checks compute, from the five gauge fields above, $D_{v}I_g = -\Sch(v)$ and $D_{v}X = v$ along every frame vector and the preservation of the flat product; the own-ray classes $(D_{\ell}I_g)\cdot\ell$ on five rational null vectors, zero on the Einstein fields (where $\Sch(a) = \tfrac{\Lambda}{6}a$ and $\mathcal{R} = 4\Lambda$), not all zero on the FLRW instant, and equal on every field to $\tfrac12\,\ell\cdot\tilde T(\ell)$ with $\tilde T(a) := -(\mathcal{G}(a) + \Lambda a)$, which is $8\pi G\,T$ of clause (7); and, in $\Cl(3,4)$ with $I_{7}$ central of square $-1$, the isotropy of $W = \mathrm{span}(I_g, \ell, s_{1} + I_{7}s_{2})$ for three null directions and the classes $(D_{\ell}I_g)\cdot w$ on it, $-\ell\cdot\Sch(s_{1}) - I_{7}\,\ell\cdot\Sch(s_{2})$ on $m$, both zero on de Sitter and the own-ray one nonzero on the FLRW instant.

\emph{In components} (a labelled reading). In the splitting of $g$ the derivatives of the setting are $\nabla_{v}(\sigma, \mu, \rho) = (\nabla_{v}\sigma - \mu(v),\ \nabla_{v}\mu + v^{\flat}\rho - \Sch(v, \cdot)\sigma,\ \nabla_{v}\rho + \Sch(v, \mu^{\sharp}))$ (Fact \claimref{interior.imported-tractor-connection}), and the clauses are the same linear algebra on the index components $\Sch_{ab}$, $\Ric_{ab}$, $T_{ab}$, $g^{ab}T_{ab} = \tau$ and $\nabla^{a}\Ric_{ab} = \tfrac12\nabla_{b}R$. The earlier checks verify that reading with exact Gaussian rationals standing for $\mathrm{span}\{1, I_{7}\}$: the class formulas of (2) and (3) on a generic rational $\Sch$, the own-ray system of (4) on eleven rational null vectors (rank nine on the ten-dimensional space of symmetric forms, kernel spanned by $g$, unchanged when the $(I_g, m)$ equations are added) and the transverse classes $-\tfrac{\Lambda}{6}g(v, w)$ at $\Lambda = 3$, the one-vacuum kernel of (6) (spanned by $\ell^{\flat}\otimes\ell^{\flat}$, trace-free), the Fock criterion of (2) in the exterior model of the complexified $\Cl(2,4)$ (for a generic, a radiation and an Einstein $\Sch$), the four-dimensional kernels of $c(I_g)$ and $c(X)$ with zero intersection on the eight-dimensional module, the real-direction count of (1) on sample planes, and for (7), with $\kappa := 4\pi G$, the store-convention $\Ric = (\Lambda + \kappa\tau)g - 2\kappa T$ with own-ray classes $\kappa T(\ell, \ell)$, trace $4\Lambda + 2\kappa\tau$, value $\Lambda g$ at $T = 0$, and the linear system of the Bianchi step forcing $dh = \kappa\,d\tau$ (and, at $T = 0$, $dR = 0$).
